\chapter{Introduction}

\section{Background and Rationale}

In recent years,  the rapid growth of technology and the push for digital 
transformation in education have made the development of robust 
Learning Management Systems (LMS) more urgent than ever. The architecture 
behind such systems is not merely a technical choice; it directly shapes 
their scalability, maintainability, and long-term. Approaches such as 
Monolithic, Modular Monolith, and Microservices have emerged as dominant 
paradigms, each offering unique strengths and trade-offs.This essay will 
explore and contrast these three architectural models in the context of 
building an LMS, focusing on functional modules such as user authentication, 
course management, content delivery, assessment, and progress tracking.
% \cite{nguyen_van_dinh_2023}.

\section{Project Objectives}

In general , the objectives of this project is building and 
deploying a Learning Management System (LMS) in form of three 
different software architectures.

For specific, the following targets will be expected:
\begin{enumerate}
    \item Design an LMS system with core functions (CRUD, 
    authentication / authorization, enrollment, 
    learning content, assignment \& grade, learning process,..)

    \item Deploy the LMS system according to three different architectures
    \item Develop scenarios for system change and expansion
    \item Compare architectures based on specific technical criteria
    \item Analyze the advantages and disadvantages of each 
    architecture in actual implementation

\end{enumerate}

\section{Scope of the Project}

The scope of the project focuses on researching, designing and comparing three 
software architectures when applied to the same business problem. The LMS system 
is only used as a research object to illustrate and verify the above architectures, 
not the main development goal of the project. In other words, the focus of the 
topic is software architecture analysis, while LMS construction only serves as 
a means for deployment and evaluation.

Across all three architectures, the system is consistently decomposed into core 
functional modules (e.g., authentication/authorization, enrollment, learning 
content, assignments and grading, progress tracking). Extended features such as 
learning analytics or recommendation algorithms are acknowledged but reserved 
for future work.

The project does not prioritize detailed user interface or user experience (UI/UX) 
design, employing only minimal interfaces sufficient for testing and demonstration. 
Regarding technology choices, the team adopts contemporary frameworks considered 
most suitable at present, namely NestJS and Vue.js due to their widespread adoption, 
strong community support, and suitability for modular.

Other architectures (e.g., serverless, large-scale event-driven architecture, 
or complex hybrid architectures) are outside the scope of the study.


\section{Research Methodology}
\lipsum[6-7]

\subsection{Theorical Methodology}
\lipsum[8-9]
\subsection{Practical Methodology}
\lipsum[9-10]

\section{Structure of the Report}
This report is structured as follows:
\begin{itemize}
    \item \textbf{Chapter 1}: Introduction - Provides background, objectives, scope, methodology, and references.
    \item \textbf{Chapter 2}: Literature Review - Reviews existing technologies and related work.
    \item \textbf{Chapter 3}: System Design and Implementation - Details the design and implementation of the messaging system.
    \item \textbf{Chapter 4}: Testing and Evaluation - Presents the testing procedures and evaluation results.
    \item \textbf{Chapter 5}: Conclusion and Future Work - Summarizes findings and suggests future directions.
    \item \textbf{Appendices}: Contains supplementary material and code listings.
\end{itemize}
\lipsum[11-14]